\SourceLeaf{482}
\Needspace{6\baselineskip}\begin{center}\bfseries SCHLUSSBEMERKUNG ZUR NEUEN AUSGABE.\end{center}

Dieser erste Band von \textsc{Gauss}' Werken ist ein Wiederabdruck der sieben Abschnitte der Disquis. Arithm., die 1801 im Oktavformat erschienen. Der achte Abschnitt, auf den an mehreren Stellen verwiesen wird und den \textsc{Gauss} wohl anfangs mit den übrigen veröffentlichen wollte, findet sich unter seinen Handschriften. Da er dessen Ausarbeitung aber nicht in gleicher Weise abgeschlossen hat wie die der ersten sieben Abschnitte, wird er sich in dieser Ausgabe den arithmetischen Abhandlungen des Nachlasses anschließen.

Textänderungen wurden in den Disqu. Ar. nur an den folgenden Stellen vorgenommen:

In Artikel 125 wurden die beiden Einschübe (wenn \(>5\)) und (\(>17\); aber \(-13N3,\allowbreak\ -17N5\)) eingefügt.

In Artikel 126. \textit{Beweis.} wurde ‘in der Reihe (I) sind \(a\) Glieder durch \(p\) teilbar, \(b\) Glieder durch \(p^2\), \(c\) Glieder durch \(p^3\) usw.’ an die Stelle von ‘in der Reihe (I) sind \(a\) Glieder durch \(p\), aber nicht durch \(p^2\) teilbar, \(b\) Glieder durch \(p^2\), aber nicht durch \(p^3\), \(c\) Glieder durch \(p^3\), aber nicht durch \(p^4\) usw.’ gesetzt.

In Artikel 128 III folgten auf ‘Denn wenn \(+b\) oder \(-b\equiv r\ (\mathrm{mod.}\ p)\) ist, so ist \(bb\equiv rr\ (\mathrm{mod.}\ 2p)\), und daher ist das Glied \(\frac12(a-bb)\) durch \(p\) teilbar.’ in der Ausgabe von 1801 noch die Worte ‘und noch viel mehr \(2(a-bb)\).’; diese wurden weggelassen.

In Artikel 129 wurde überall \(2\surd a+1\) an die Stelle von \(2\surd a\) und demnach 9 an die Stelle von 4 als diejenige Zahl gesetzt, unter die \(a\) jedenfalls nicht sinkt.

In Artikel 139 lautete der Schluss der Untersuchung über den Fall (4): ‘Man erkennt leicht, dass sich aus dieser Gleichung \(a'pRh,\allowbreak\ \pm ahpRa',\allowbreak\ \pm aa'hRp\) herleiten lässt; dies stimmt mit dem in (2) Gefundenen überein. In den übrigen Fällen ist der Beweis derselbe.’ Dies wurde entsprechend der Anmerkung geändert, die \textsc{Gauss} dem Artikel 2 der Abhandlung \textit{Theorematis arithmetici demonstratio nova} beigefügt hat: ‘Es wird nicht unangebracht sein, hier einen kleinen Fehler anzugeben und zu berichtigen, der sich durch eine mir unbekannte Nachlässigkeit in die Darstellung jenes Beweises eingeschlichen hat. Auf Seite (108) sind von Zeile (14) an die folgenden Überlegungen einzusetzen: Man erkennt leicht, dass sich aus dieser Gleichung \(a'pRh\ldots(\alpha);\allowbreak\ \pm ahRa'\ldots(\delta);\allowbreak\ \pm ahRp\ldots(\gamma)\) herleiten lässt. Aus \((\alpha)\), das mit \((\alpha)\) in (2) übereinstimmt, folgt ebenso wie dort, dass entweder zugleich \(hRp,\allowbreak\ hRa'\) oder \(hNp,\allowbreak\ hNa'\) gilt. Im ersten Fall würde aber aus \((\delta)\) folgen, dass \(aRa'\) gilt, entgegen der Voraussetzung; daher ist \(hNp\) und somit nach \((\gamma)\) auch \(aNp\).

\SourceLeaf{483}
In Artikel 171 wurde nach einer handschriftlichen Bemerkung ‘in der \(A\) weder größer als \(\surd{\frac43D},\allowbreak\ C\), noch kleiner als \(2B\) ist’ an die Stelle von ‘in der \(A\) nicht \(>\surd{\frac43D}\), \(B\) nicht \(>\frac12A\), \(C\) nicht \(<A\) ist’ gesetzt.

In Artikel 174. \textit{Zweite Methode.} wurde ‘Da \(a\) nicht \(>\surd{\frac43D}\) sein wird, werden offenbar alle auf diese Weise entstehenden Formen reduziert sein.’ an die Stelle von ‘Entstehen auf diese Weise Formen, in denen \(a>\surd{\frac43D}\) ist, so sind sie zu verwerfen; alle übrigen werden dagegen offenbar reduziert sein.’ gesetzt.

\Needspace{6\baselineskip}In Artikel 256 VI enthält die frühere Ausgabe unter den positiven Determinanten der Form \(8n+5\), für die die Anzahl der Klassen in der eigentlich primitiven Ordnung dreimal so groß ist wie die in der uneigentlich primitiven, auch die Determinante 397; zu ihr gehören jedoch gleich viele Klassen beider Arten.

In Artikel 302 wird nach einer handschriftlichen Bemerkung die Anzahl der Klassen für das zweite Tausend negativer Determinanten mit 28595 angegeben und nicht mit 28603 wie in der früheren Ausgabe.

In Artikel 303 werden die Anzahlen der Determinanten für die Klassifikationen II. 4; II. 5; IV. 2 nach den im Nachlass gefundenen Tafeln mit 31, 44, 69 angegeben und nicht mit 32, 42, 68 wie in der früheren Ausgabe.

In Artikel 325 wurde ‘\(23t,\allowbreak\ 23t\pm1,\allowbreak\ 23t\pm5,\allowbreak\ 23t\pm7,\allowbreak\ 23t\pm9,\allowbreak\ 23t\pm10\). Nach deren Streichung bleiben 119, 127 übrig, die allein \(V\) einen quadratischen Wert geben’ an die Stelle von ‘\(23t,\allowbreak\ 23t\pm5,\allowbreak\ 23t\pm7\), \(23t\pm9,\allowbreak\ 23t\pm10\). Nach deren Streichung bleiben 119, 127, 137 übrig, von denen nur die ersten beiden \(V\) einen quadratischen Wert geben.’ gesetzt.

In Artikel 360 IV wurden die Worte ‘während bei den meisten anderen kubischen Gleichungen, deren Wurzeln sämtlich reell sind, die Dreiteilung sowohl eines Winkels als auch eines Verhältnisses nicht vermieden werden kann.’ weggelassen, die auf ‘zum Beispiel für \(\beta=3\) allein die Dreiteilung eines Winkels’ folgten und deren Unrichtigkeit \textsc{Gauss} in seinem Handexemplar vermerkt hat.

Die Anmerkungen auf den Seiten 80, 102, 133, 211, 263, 265, 449 sind dem handschriftlichen Nachlass entnommen.

Um die Übersicht zu erleichtern, wurde die äußere Form der neuen Ausgabe gegenüber dem Druck dieses Werkes von 1801 in einigen Punkten verändert. Dazu glaubte man sich umso mehr berechtigt, als \textsc{Gauss} nach eigener Aussage auch in anderen Punkten auf Platzersparnis Rücksicht nahm. Viele Formeln, die in der Textzeile standen, wurden gesondert herausgesetzt. Die Inhaltsangaben, die gesammelt vor dem Werk standen, wurden auch den einzelnen Abteilungen vorangestellt und nicht nur den Abschnitten wie in der ersten Ausgabe; die allgemeineren darunter wurden als Seitenüberschriften verwendet.

\BookRule{45mm}

\begin{center}\small GÖTTINGEN,\\ GEDRUCKT IN DER DIETERICHSCHEN UNIVERSITÄTS-DRUCKEREI\\ W. FR. KAESTNER.\end{center}
